\section{The minimal reversible localizer and the proposed simple types}

\begin{claim}{D.1}{Proved}{Minimality of integral homology: the valid part of Lead 3}
The class $\WZ$ of all functors inducing integral homology isomorphisms on nerves is a fundamental localizer satisfying Loc (6 a,b,c). It is the smallest modern fundamental localizer satisfying Loc (6 b), and is strictly larger than $\Winfty$.
\end{claim}
\begin{proof}
Identities, two out of three, and the weak retraction axiom follow by applying the functor $H_*(-;\mathbb Z)$. A terminal object contracts the nerve. To verify relative Theorem A without presupposing that this class is fundamental, take a triangle $A\xrightarrow{u}B\to C$ with each $A/c\to B/c$ a homology equivalence. The augmentations
\[
 \int_{c\in C}(A/c)\longrightarrow A,\qquad
 \int_{c\in C}(B/c)\longrightarrow B
\]
are weak homotopy equivalences. For example the first sends $(c,a,\gamma:v u(a)\to c)$ to $a$; it is right adjoint to the section $a\mapsto(vu(a),a,\id)$, since a morphism out of this section is uniquely determined by the arrow in $A$. The augmentations commute with $u$. The simplicial replacement of either diagram has the first-quadrant homology spectral sequence
\[
 E^1_{p,q}=\bigoplus_{c_0\to\cdots\to c_p}H_q(B(A/c_0);\mathbb Z)
 \ \Longrightarrow\ H_{p+q}\left(B\int_C(A/c);\mathbb Z\right),
\]
and likewise for $B$. This is the normalized double chain complex of the replacement; its filtration is finite in each total degree. The induced map is an isomorphism on the $E^1$ page, hence on homology. The two augmentations prove that $u\in\WZ$.

For a sieve cover $X=X_1\cup X_2$ with intersection $X_0$, every nondegenerate chain lies in a piece containing its last vertex. Thus there is a natural short exact sequence of free chain complexes
\[
 0\to C_*(NX_0;\mathbb Z)
 \xrightarrow{(i_*,-j_*)}C_*(NX_1;\mathbb Z)\oplus C_*(NX_2;\mathbb Z)
 \to C_*(NX;\mathbb Z)\to0.
\]
For a map of such squares, take the algebraic cones of the four induced chain maps. They form the same short exact sequence, with the middle term the direct sum of the two piece cones. The homology long exact sequence says that acyclicity of any two of the intersection cone, the middle cone, and the union cone implies acyclicity of the third. A direct sum of two complexes is acyclic exactly when each summand is acyclic. This proves all three cube implications, including every empty-piece case, using unreduced chains. The reverse-edge version follows from the degenerate cube.

Every other fundamental localizer satisfying the reverse axiom contains $\WZ$ by \claimref{B.6}. Strictness over $\Winfty$ is witnessed by the functor $A\to e$ in \claimref{B.7}. The assertion that $\Winfty$ satisfies the complete manuscript list including Loc (6 b) is therefore false, not merely unsupported by its proposed proof.
\end{proof}

\begin{claim}{D.2}{Proved}{The incompatibility of the two poles, with the connectivity correction}
There is no modern fundamental localizer satisfying Loc (5), Loc ($7\mathrm D_{\Winfty}$), and Loc (6 b). If Loc (5) is omitted, the unique fundamental localizer satisfying both Loc ($7\mathrm D_{\Winfty}$) and Loc (6 b) is $\Fl(\Cat)$.
\end{claim}
\begin{proof}
The fibration classification proved in Part E says that Loc ($7\mathrm D_{\Winfty}$) leaves exactly $\Winfty$ and the total class. The first fails Loc (6 b) by \claimref{B.8}. The total class satisfies every stability implication, including the two in this statement, but it fails Loc (5), since $\varnothing\to e$ has a nonbijective component map. This last point corrects the wording of Lead 3: the total class is the answer only after dropping the connectedness requirement. It is not an exceptional solution satisfying that requirement.
\end{proof}

\begin{claim}{D.3}{Proved}{A nontotal reversible localizer kills no sphere}
If a fundamental localizer satisfies Loc (6 b) and makes any sphere $S^n$, $n\geq0$, aspherical, it is the total class. In particular no nontotal reversible localizer contains $\W_1$.
\end{claim}
\begin{proof}
The cofiber criterion applied to $X\to e$ identifies asphericity of $X$ with asphericity of its unreduced suspension, for every category $X$. The categorical models of spheres supplied by simplex categories permit repeated application. Starting with $S^n$, we obtain asphericity of $S^0$. But $K(\varnothing\to e)=e\amalg e$ represents $S^0$; one more application gives $\varnothing\to e\in\W$.

For clarity, this forces totality under the stated axioms. Fundamental localizers are stable under product with any category $C$: for $u:A\to B$ in $\W$, the slices over $c\in C$ of $u\times C$ identify with $u\times(C/c)$; the projections to $u$ are homotopy equivalences because $C/c$ has a terminal object. Relative Theorem A gives $u\times C\in\W$. Applying this to $\varnothing\to e$ gives $\varnothing\to C\in\W$ for every $C$. Two out of three on $\varnothing\to A\to B$ gives every functor $A\to B$. Finally $\W_1$ contains $S^2\to e$ by its defining homotopy-group test.
\end{proof}

\begin{claim}{D.4}{Proved}{Both readings of the stability intersection: Lead 4}
Let $\widetilde\W$ be the intersection of all modern fundamental localizers satisfying the stability list XVI Loc (1 ter, 2, 3 bis, 4, 6, 7, 8), with no prior inclusion in $\Winfty$ imposed. If Loc (6) includes its reverse half, then $\widetilde\W=\Fl(\Cat)$. If it means only Loc (6 a), then $\widetilde\W=\Winfty$. If the definition additionally insists, as the prose of XVI ms.~87 does, that $\widetilde\W\subseteq\Winfty$, there is no such localizer in the first reading.
\end{claim}
\begin{proof}
Loc (7) implies Loc ($7\mathrm D_{\Winfty}$), as made precise in Parts A and E. With the reverse half of Loc (6), \claimref{D.2} forces each member of the intersection to be total. The total class itself satisfies the entire stability list, so the intersection is total. This class is not a subclass of $\Winfty$, proving the last sentence.

For the forward reading, minimality \claimref{B.1} gives the lower bound $\Winfty$. This class satisfies strong saturation, the terminal-object and relative Theorem A axioms, disjoint-union preservation and reflection, forward gluing, the ordinary homotopy-fibration axiom, and filtered-colimit stability. The first six assertions are established in Parts A--C; filtered-colimit stability follows because each finite simplicial homotopy problem factors through a stage of a filtered diagram of nerves. Thus $\Winfty$ belongs to the family being intersected, giving the opposite inclusion. No conjecture about simple homotopy can yield a strict intermediate answer under either reading.
\end{proof}

\begin{claim}{D.5}{Proved}{What can and cannot retain simple homotopy information}
No modern fundamental localizer on all small categories is a strict subclass of $\Winfty$. Even weakening relative Theorem A to the absolute version while retaining weak saturation and terminal asphericity does not produce such a subclass. A finer equivalence relation can instead be retained on finite CW or finite simplicial complexes by keeping elementary expansions, collapses, and the simple structures as geometric data. Such a structure is not a fundamental localizer. Merely localizing the category of all finite CW complexes and all cellular maps at simple homotopy equivalences does not retain the simple homotopy classification.
\end{claim}
\begin{proof}
The first assertion is \claimref{B.1}, and the stronger weak-localizer minimum is \claimref{A.19}. For the last assertion, the endpoint inclusions $X\to X\times[0,1]$ of a finite CW complex are simple homotopy equivalences: collapse the prisms one cell at a time. Hence localization at simple homotopy equivalences identifies homotopic cellular maps. If $f:X\to Y$ is any cellular homotopy equivalence with cellular homotopy inverse $g$, then $gf$ and $fg$ become identities there, making $f$ invertible whether or not it is simple. Thus a homotopy-category localization of this unrestricted mapping category loses the distinction in question.

There is also an immediate obstruction at the saturation axiom. Whenever $f$ is a nonsimple homotopy equivalence and $g$ a homotopy inverse, both $gf$ and $fg$ are homotopic to identities and are simple, while $f$ is not. Thus the printed Loc (1 c), and a fortiori the fundamental-localizer package, cannot hold for the class of simple equivalences in this mapping category. This is a conditional obstruction for every such $f$, not a claim that Whitehead torsion is an invariant of an undecorated homotopy type.

A suitable weaker structure is the groupoid generated by elementary expansions, elementary collapses, and cellular isomorphisms, retaining the allowed simple structures and the relations one chooses between these geometric operations. Its components classify simple homotopy types by definition. Alternatively one retains torsion in a category of structured finite complexes, instead of declaring all homotopy inverses invertible through the fundamental-localizer axioms. The classical definitions and the based chain construction of the torsion lie in \citet[Definition 2, Example 3, and Definition 17]{RootLurie2014Whitehead}; the Whitehead group is
\[
 \operatorname{Wh}(\pi)=K_1(\mathbb Z[\pi])/\langle[\pm g]:g\in\pi\rangle.
\]
It is not merely an unspecified cohomological invariant in a group denoted $K(\mathbb Z[\pi_1])$. This distinguishes the valid geometric suggestion in XVI ms.~88--96 from the impossible proposed fundamental localizer.
\end{proof}

\begin{claim}{D.6}{Proved}{The finite-poset observation survives without the conjecture}
Let $P$ be the face poset of the boundary of a triangle, consisting of its three vertices and three edges, and let $\operatorname{sd}P$ be the poset of nonempty chains in $P$. The last-element map $\operatorname{sd}P\to P$ is in $\Winfty$ but has no inverse up to zigzags of natural transformations of functors. Thus strict categorical homotopy is finer than weak homotopy of nerves; this fact does not produce a fundamental localizer smaller than $\Winfty$.
\end{claim}
\begin{proof}
The two nerves realize to circles. The canonical subdivision homeomorphism and the usual linear homotopy within each simplex identify realization of the last-element map with a degree-one map. It is consequently a weak homotopy equivalence. The nerve of $P$ is a cycle with six edges, while that of $\operatorname{sd}P$ is a cycle with twelve edges. Any order-preserving map $P\to\operatorname{sd}P$ induces a simplicial map taking each of the six edges to one edge or a vertex. A map of nonzero degree between these oriented circles must traverse every one of the twelve target edges, so no such simplicial map has nonzero degree. A natural transformation gives an actual homotopy on nerves, as does any finite zigzag. Thus no order-preserving map in the reverse direction can be a categorical homotopy inverse of the last-element map. The program \texttt{code/poset\_subdivision.py} enumerates every order-preserving map and verifies the vanishing degree; its data and output are attached. This checks the relevant observation at XVI ms.~95 while leaving minimality intact.
\end{proof}
